\documentclass{baustms}
\citesort
\theoremstyle{cupthm}
\newtheorem{theorem}{Theorem}[section]
\newtheorem{lemma}[theorem]{Lemma}
\newtheorem{proposition}[theorem]{Proposition}
\theoremstyle{cupdefn}
\newtheorem{definition}[theorem]{Definition}
\newtheorem{convention}[theorem]{Convention}
\newtheorem{example}[theorem]{Example}
\theoremstyle{cuprem}
\newtheorem{remark}[theorem]{Remark}

\begin{document}
\runningtitle{Two degrees or two face sizes}
\title{No alternating plane graph has exactly two vertex degrees,
       or exactly two face sizes}
\author[1]{Hanyu Yang}
\address[1]{Independent researcher, California, USA\email{hanyu.yang.92@gmail.com}}
\authorheadline{H. Yang}
\classification{primary 05C10; secondary 05C30}
\keywords{alternating plane graph, plane graph, discharging, Euler's formula}

\begin{abstract}
A plane graph is \emph{alternating} if every vertex has degree at least three,
every face has size at least three, adjacent vertices have different degrees and
adjacent faces have different sizes. Alth\"ofer, Haugland, Scherer, Schneider and
Van Cleemput introduced these graphs and conjectured that none has exactly two
degrees, and none exactly two face sizes. We prove both, by one per-edge
inequality applied twice with the roles of vertices and faces exchanged. Summing the reciprocals of the two endpoint degrees and the two
incident face sizes over all edges returns the number of vertices plus the number
of faces; parity and alternation then force that sum to be at most one on every
edge, while Euler's formula forces some edge past one. We also prove both
statements when face alternation is required only between distinct adjacent
faces, so that bridges are admitted; that extension uses a charging argument.
\end{abstract}

\section{Introduction}

Alth\"ofer, Haugland, Scherer, Schneider and Van Cleemput~\cite{AHSSV2015}
introduced alternating plane graphs: plane graphs in which adjacent vertices
have different degrees and adjacent faces have different sizes. The condition is
severe, and it makes even existence questions hard. Their paper closes
(\cite{AHSSV2015}, \S10, p.~362) with four open problems, of which the first is:

\begin{quote}
\textbf{Conjecture 10.1.} There are no $2,Y$-alternating plane graphs and no
$X,2$-alternating plane graphs.
\end{quote}

Here an $X,Y$-alternating plane graph has exactly $X$ distinct vertex degrees
and exactly $Y$ distinct face sizes. We prove the conjecture
(Theorem~\ref{thm:101}); equivalently, every alternating plane graph has at least
three distinct vertex degrees and at least three distinct face sizes.

What was known is this. Section~3.2 of \cite{AHSSV2015} shows that a
$2,Y$-alternating plane graph has degrees $3$ and $4$, or degrees $3$ and $5$,
and bounds its order below by $25$ in the first case and by $56$ in the second;
the search of its \S4 reaches neither. Separately, its \S9.1 bounds the number of
faces of a \emph{weak} alternating plane graph --- one in which vertices of
degree $2$ are permitted --- with degrees $2$ and $k$, by summing $1/r+1/s$ over
the edges separating an $r$-face from an $s$-face. That face-incidence count
never uses the degree $2$. Combining it with the analogous count over vertices
excludes both surviving degree pairs at once, and, applied with the roles of
vertices and faces exchanged, excludes two face sizes as well.

Why some such combination is needed can be seen immediately. Alternation alone
gives at most $\frac13+\frac14$ from the two endpoint degrees of an edge and at
most $\frac13+\frac14$ from its two incident faces, hence only
$\frac{7}{6}$ in total, and $\frac76 > 1$ is compatible with Euler's formula. The
two-valued hypothesis supplies what is missing: it makes the graph bipartite on
the side where the two values sit, so on the \emph{other} side the two values are
even, hence at least $4$ and $6$, and that contribution drops to
$\frac14+\frac16=\frac5{12}$. The two sides then total exactly $1$, and Euler's
formula is contradicted. Both halves of the conjecture are this one argument.

The note is organised so that its two achievements stay separate.
Section~\ref{sec:prelim} fixes the definitions and the two conventions that the
source paper leaves implicit. Section~\ref{sec:identity} isolates the edge charge
and the local parity lemma, with a worked example. Section~\ref{sec:101} proves
Conjecture~10.1 as stated. Section~\ref{sec:c2} then proves an extension, not a
prerequisite: both nonexistence statements survive when face alternation is
required only between \emph{distinct} adjacent faces, so that bridges are
admitted. That extension is the longer and harder argument, and its bottleneck is
identified there.

\section{Preliminaries}\label{sec:prelim}

We follow \cite{AHSSV2015} throughout. All graphs are finite, simple and plane;
those in \S\ref{sec:identity}--\S\ref{sec:c2} are also connected. Faces include
the unbounded one.

\begin{definition}[\cite{AHSSV2015}, Definition~2.1]\label{def:apg}
An \emph{alternating plane graph} is a connected plane graph in which
\begin{enumerate}
  \item[(a)] every vertex has degree at least $3$;
  \item[(b)] every face has size at least $3$;
  \item[(c)] adjacent vertices have different degrees;
  \item[(d)] adjacent faces have different sizes.
\end{enumerate}
\end{definition}

\begin{definition}[\cite{AHSSV2015}, Definitions~3.1 and~3.4]
A $(3,4,5)$-alternating plane graph is one in which every vertex degree and
every face size lies in $\{3,4,5\}$. An \emph{$X,Y$-alternating plane graph} is
one with exactly $X$ distinct vertex degrees and exactly $Y$ distinct face
sizes.
\end{definition}

\begin{convention}[face size]\label{conv:c1}
The \emph{size} of a face is the length of its boundary walk, that is, the number
of edge-side incidences it carries. In particular a bridge, lying on one face
from both sides, contributes $2$ to that face's size.

Boundary-walk length is not the number of distinct vertices on the boundary. Two
quadrilaterals meeting in exactly one vertex have an exterior boundary walk of
length $8$, although only $7$ distinct vertices lie on it. (This illustrates the
convention; it is not an alternating plane graph.) The distinction matters
because it is boundary-walk length, not the vertex count, that is forced to be
even in a bipartite plane graph.

\cite{AHSSV2015} does not define ``size'' separately, but its arguments use
edge-side incidences: this is the count underlying its equation~(3.1) and its
face-incidence identity~(9.2), and it is what makes the assertion in the proof of
its Lemma~9.2 that a bipartite graph has $f_j = 0$ for every odd $j$ correct.
Page~339 also treats a bridge as producing a face adjacent to itself, so the
paper's face language is not confined to simple cycles.
\end{convention}

\begin{convention}[bridges]\label{conv:c2}
In Definition~\ref{def:apg}, two face occurrences on opposite sides of an edge
must have different sizes. A bridge has one face on both sides, so bridges are
excluded; \cite{AHSSV2015} says so on p.~339, where an alternating plane graph is
noted to be ``always at least $2$-edge-connected, since plane graph with edge
connectivity $1$ contains a face that is adjacent to itself''.

Call a plane graph \emph{bridge-relaxed alternating} if it satisfies
Definition~\ref{def:apg} with (d) imposed only on pairs of \emph{distinct}
adjacent faces. Bridges are then permitted. Face adjacency is always meant
through opposite sides of an edge, never through a shared vertex.
\end{convention}

Note that \cite{AHSSV2015}'s ``weak'' alternating plane graphs relax a different
condition, namely (a), by admitting degree~$2$. The two relaxations are
independent, and we use ``weak'' only in the source's sense.

\section{Edge charge and local parity}\label{sec:identity}

We isolate the face-incidence count used in \cite{AHSSV2015}, \S9.1, and combine
it with the analogous count over vertices.

\begin{lemma}\label{lem:identity}
Let $G$ be a finite connected simple plane graph with at least one edge, with $v$
vertices, $e$ edges and $f$ faces. For an edge $a$ with ends $u,w$ and face
occurrences $F^-(a)$, $F^+(a)$ on its two sides, put
\[
  c(a) \;=\; \frac{1}{\deg u} + \frac{1}{\deg w}
           + \frac{1}{|F^-(a)|} + \frac{1}{|F^+(a)|}.
\]
Then, by Euler's formula,
\begin{equation}\label{eq:identity}
  \sum_{a} c(a) \;=\; v + f \;=\; e + 2,
  \qquad\text{equivalently}\qquad
  \sum_{a} \bigl(c(a) - 1\bigr) \;=\; 2 .
\end{equation}
\end{lemma}

\begin{proof}
A vertex of degree $d$ is an end of exactly $d$ edges and contributes $1/d$ each
time, hence $1$ in total; summing gives $v$. By Convention~\ref{conv:c1} a face
of size $s$ occurs exactly $s$ times as a side of an edge and contributes $1/s$
each time, hence $1$ in total; summing gives $f$.
\end{proof}

Neither count requires facial walks to be simple: repeated vertices, cut vertices
and bridges are all permitted, a bridge simply having $F^-(a) = F^+(a)$ and
contributing $2/|F|$ from its face side. This is what makes
Lemma~\ref{lem:identity} usable in \S\ref{sec:c2}, where bridges are the object
of study. We write $x(a) = c(a) - 1$ for the \emph{excess} of an edge, so
that~\eqref{eq:identity} reads $\sum_a x(a) = 2$: some edge must have positive
excess.

\begin{example}\label{ex:worked}
Start with a triangle and replace each of its three sides by two internally
disjoint paths of length two, drawn alongside that side. The result is a simple
plane graph with $9$ vertices, $12$ edges and $5$ faces: three quadrilaterals,
one for each original side, together with an inner and an outer hexagon. The
three original corners have degree $4$ and the six new vertices degree $2$, so
every edge joins a degree-$2$ vertex to a degree-$4$ vertex; and every edge lies
on one quadrilateral and one hexagon. Hence every edge has
\[
  c(a) \;=\; \tfrac12 + \tfrac14 + \tfrac14 + \tfrac16 \;=\; \tfrac76,
  \qquad
  \sum_a x(a) \;=\; 12 \cdot \tfrac16 \;=\; 2 ,
\]
as Lemma~\ref{lem:identity} requires. This graph satisfies both alternation
conditions and fails to be an alternating plane graph only through its minimum
degree. It also attains $f = \tfrac{5}{12}e$ with $f=5$ and $e=12$, so the face
bound $f \le \tfrac5{12}e$ of (\cite{AHSSV2015}, (9.5)) is sharp and cannot be
improved to any $\tfrac5{12}-\varepsilon$.
\end{example}

\begin{remark}\label{rem:calibration}
Run on the weak alternating plane graphs with degrees $2$ and $k$, $k \ge 3$,
that \cite{AHSSV2015} treats, the same estimate gives
$c(a) \le \tfrac12+\tfrac1k+\tfrac5{12} = \tfrac{11}{12}+\tfrac1k$, which is at
most $1$ exactly when $k \ge 12$; \cite{AHSSV2015} excludes $k = 11$ by a
separate argument (its Lemma~9.4) and exhibits such graphs for $k \le 10$. This
shows where the hypothesis $\deg \ge 3$ enters below. Agreement of the threshold
is not an independent verification of the face bound; Example~\ref{ex:worked},
which attains the bound, is the check that would detect a strengthening of it.
\end{remark}

The second local ingredient is a parity statement. Stating it over the corners at
a vertex, rather than over ``the incident faces'', is what makes it correct at a
cut vertex, where one face occupies several corners: those corners are
non-consecutive, so the alternation below is undisturbed.

\begin{lemma}[parity]\label{lem:parity}
Let $G$ be a plane graph whose faces admit a colouring by two colours that is
proper on every pair of distinct adjacent faces. Then for every vertex $u$,
\[
  \deg u \;-\; \bigl|\{\text{bridges at } u\}\bigr| \quad\text{is even.}
\]
\end{lemma}

\begin{proof}
List the edges at $u$ in cyclic order as $a_1,\dots,a_d$, $d = \deg u$, and let
$F_i$ be the face occurrence in the corner between $a_i$ and $a_{i+1}$, indices
modulo $d$, so that $F_{i-1}$ and $F_i$ are the two sides of $a_i$. These two
occurrences belong to the same face precisely when $a_i$ is a bridge. Hence the
corner colour changes across each non-bridge edge at $u$ and is unchanged across
each bridge at $u$. The cyclic list closes up, so the number of changes is even,
and that number is $\deg u$ minus the number of bridges at $u$.
\end{proof}

In particular, if no bridge is present at $u$ then $\deg u$ is even. What the
lemma controls is degree minus bridge incidence, not degree itself: in a triangle
with a pendant edge the four degrees are $3,2,2,1$, the numbers of incident
bridges are $1,0,0,1$, and the differences $2,2,2,0$ are even, although no vertex
has even degree. (That graph illustrates the lemma, not
Definition~\ref{def:apg}.) In an $X,2$-alternating plane graph, and in a
bridge-relaxed one, the face size supplies the required colouring.

\section{Conjecture 10.1}\label{sec:101}

\begin{theorem}\label{thm:101}
There is no $2,Y$-alternating plane graph and no $X,2$-alternating plane graph.
\end{theorem}

The two cases are the same inequality with the roles of vertices and faces
exchanged, but they are not formally dual and we prove both. Passing to the dual
graph would need $3$-edge-connectivity (\cite{AHSSV2015}, Lemma~2.2): a
$2$-edge cut produces parallel edges in the dual, which therefore need not lie in
the simple-graph class we work in. We exchange the two local counting arguments
instead.

\begin{proof}
\emph{Case 1: exactly two vertex degrees $3 \le d_1 < d_2$.}
Adjacent vertices differ in degree and only two degrees occur, so every edge
joins a $d_1$-vertex to a $d_2$-vertex. In particular $G$ is bipartite, and every
edge contributes
\[
  \frac{1}{d_1} + \frac{1}{d_2} \;\le\; \tfrac13 + \tfrac14 \;=\; \tfrac{7}{12}
\]
to the vertex side, using $d_1 \ge 3$ from Definition~\ref{def:apg}(a).

Since $G$ is bipartite, every closed walk has even length, so by
Convention~\ref{conv:c1} every face has even size, hence size at least $4$. By
Convention~\ref{conv:c2} the two sides of an edge are distinct faces, and by
Definition~\ref{def:apg}(d) their sizes differ; being distinct even numbers at
least $4$, one is at least $4$ and the other at least $6$, so the face side
contributes at most $\tfrac14 + \tfrac16 = \tfrac{5}{12}$.

Hence $c(a) \le \tfrac{7}{12} + \tfrac{5}{12} = 1$ for every edge, so
$\sum_a x(a) \le 0$, contradicting~\eqref{eq:identity}.

\smallskip
\emph{Case 2: exactly two face sizes $3 \le s_1 < s_2$.}
By Convention~\ref{conv:c2} the two sides of every edge are distinct faces, and
by Definition~\ref{def:apg}(d) they have different sizes; with only two sizes
available, every edge separates an $s_1$-face from an $s_2$-face and contributes
at most $\tfrac13 + \tfrac14 = \tfrac{7}{12}$ to the face side.

The face sizes two-colour the faces, properly on adjacent pairs, and there are no
bridges; so by Lemma~\ref{lem:parity} every degree is even. Degrees are at least
$3$ and even, hence at least $4$; adjacent vertices have distinct even degrees,
so the vertex side contributes at most $\tfrac14+\tfrac16 = \tfrac{5}{12}$. Again
$c(a) \le 1$ for every edge and~\eqref{eq:identity} is contradicted.
\end{proof}

Case~1 improves (\cite{AHSSV2015}, (9.9)) by replacing the degree-$2$
contribution $\tfrac12$ of its (9.8) by $1/d_1 \le \tfrac13$; the face bound
(9.5) it is combined with is already there, and its derivation never used the
degree $2$.

\section{The bridge-relaxed extension}\label{sec:c2}

Definition~\ref{def:apg} excludes bridges, so Theorem~\ref{thm:101} is proved.
We now show that both conclusions survive the weaker requirement of
Convention~\ref{conv:c2}, in which face alternation is imposed only between
distinct adjacent faces and bridges are therefore admitted.

The relaxation introduces exactly one new difficulty: a bridge need not force a
change of face size, so a vertex on a bridge need not have even degree, and the
argument of Case~2 above breaks down. The following lemma restores enough
control. It says that a bridge is expensive --- its face is large --- and the
proof of Theorem~\ref{thm:x2weak} then pays for the few edges that can have
positive excess out of that expense. The essential step there, and the only place
where real work is done, is to bound the number of positive edges assigned to
each bridge.

\begin{lemma}\label{lem:bridgeface}
In a finite connected simple plane graph of minimum degree at least $3$, every
face incident with a bridge has boundary-walk length at least $8$.
\end{lemma}

\begin{proof}
Let $a = uw$ be a bridge and let $A$ and $C$ be the components of $G-a$
containing $u$ and $w$. Then $\deg_A u = \deg u - 1 \ge 2$ and likewise
$\deg_C w \ge 2$, so each of $A$ and $C$ has minimum degree at least $1$ and at
least three vertices. Let $W_A$ be the facial boundary walk of $A$ through the
corner exposed by deleting $a$. Its length is not $1$, which would need a loop.
Nor is it $2$: that is either a digon of parallel edges, excluded by simplicity,
or a single edge traversed in both directions, and for those two darts to
constitute the whole facial boundary both ends of that edge would have degree
$1$, forcing $A = K_2$ against $\deg_A u \ge 2$. So $|W_A| \ge 3$, and likewise
$|W_C| \ge 3$. The face $F$ of $G$ incident with both sides of $a$ is traversed
as $W_A$, then $a$, then $W_C$, then $a$ in the opposite direction, so by
Convention~\ref{conv:c1}
\[
  |F| \;=\; |W_A| + |W_C| + 2 \;\ge\; 3 + 3 + 2 \;=\; 8 . \qedhere
\]
\end{proof}

\begin{proposition}\label{prop:twoY-weak}
There is no bridge-relaxed alternating plane graph with exactly two vertex
degrees.
\end{proposition}

\begin{proof}
Let $G$ be one, with degrees $3 \le d_1 < d_2$. As in Case~1 of
Theorem~\ref{thm:101}, every edge joins a $d_1$-vertex to a $d_2$-vertex, so $G$
is bipartite; hence every closed walk has even length --- a facial walk
traversing a bridge twice included --- and every face has even size at least $4$.
A bridge lies on a face of size at least $8$ by Lemma~\ref{lem:bridgeface}, so it
has
\[
  c(a) \;\le\; \tfrac13 + \tfrac14 + \tfrac28 \;=\; \tfrac56 ,
\]
while a non-bridge edge, whose two sides are distinct faces of distinct even
sizes at least $4$ and $6$, has
$c(a) \le \tfrac13+\tfrac14+\tfrac14+\tfrac16 = 1$. So $x(a) \le 0$ for every
edge, contradicting~\eqref{eq:identity}.
\end{proof}

\begin{theorem}\label{thm:x2weak}
There is no bridge-relaxed alternating plane graph with exactly two face sizes.
\end{theorem}

\begin{proof}
Let $G$ be one, with face sizes $s_1 < s_2$. If $G$ has no bridge then it
satisfies Definition~\ref{def:apg} outright and is excluded by
Theorem~\ref{thm:101}. So assume $G$ has $B \ge 1$ bridges.

\smallskip
\emph{Step 1: only an edge incident with a vertex of degree $3$ can have positive
excess.} By Lemma~\ref{lem:bridgeface} every face carrying a bridge has size at
least $8$; since that face is an $s_1$- or an $s_2$-face and $s_1 < s_2$, this
gives $s_2 \ge 8$ as well. Every non-bridge edge separates an $s_1$-face from an
$s_2$-face, so its face side is at most $\tfrac1{s_1}+\tfrac1{s_2} \le
\tfrac13+\tfrac18$. Degrees at the two ends of an edge are distinct by
Definition~\ref{def:apg}(c) and at least $3$. Three cases:
\begin{itemize}
  \item $a$ a bridge, its face of size at least $8$:
        $c(a) \le \tfrac13+\tfrac14+\tfrac28 = \tfrac56$, so
        $x(a) \le -\tfrac16$;
  \item $a$ not a bridge, both ends of degree at least $4$:
        $c(a) \le \tfrac14+\tfrac15+\tfrac13+\tfrac18 = \tfrac{109}{120}$, so
        $x(a) \le -\tfrac{11}{120}$;
  \item $a$ not a bridge, one end of degree $3$:
        $c(a) \le \tfrac13+\tfrac14+\tfrac13+\tfrac18 = \tfrac{25}{24}$, so
        $x(a) \le \tfrac1{24}$.
\end{itemize}
Only the third class can be positive. Call its members \emph{positive-eligible}
and let $P$ be their number.

\smallskip
\emph{Step 2: each bridge receives at most two positive-eligible edges.} This is
the bottleneck of the proof. Let $a$ be positive-eligible and let $u$ be its end
of degree $3$, which is unique since adjacent vertices have different degrees. By
Lemma~\ref{lem:parity} --- applicable because the two face sizes colour the faces
properly on distinct adjacent pairs --- the number of bridges at $u$ has the same
parity as $\deg u = 3$, so it is $1$ or $3$; it is not $3$, since $a$ itself is a
non-bridge edge at $u$. So $u$ lies on exactly one bridge, written $b(u)$, and we
\emph{assign} $a$ to $b(u)$. The assignment is therefore defined on every
positive-eligible edge and single-valued. The non-bridge incidence is essential:
a degree-$3$ vertex with three incident bridges would receive no such assignment.

Now fix a bridge $b$. Every edge assigned to $b$ has its degree-$3$ end at a
vertex whose unique bridge is $b$, so that vertex is an end of $b$. The two ends
of $b$ have different degrees by Definition~\ref{def:apg}(c), so at most one of
them has degree $3$; call it $u$ if it exists. Of the three edges at $u$ one is
$b$ itself, leaving exactly two that can be assigned to $b$. Hence $P \le 2B$.

\smallskip
\emph{Step 3: summing the excesses.} Partition the edges into the $P$
positive-eligible ones, the $B$ bridges, and the rest, whose excesses are at most
$\tfrac1{24}$, $-\tfrac16$ and $0$ respectively by Step~1. Then
by~\eqref{eq:identity},
\begin{equation}\label{eq:final}
  2 \;=\; \sum_a x(a) \;\le\; \frac{P}{24} - \frac{B}{6}
     \;\le\; \frac{2B}{24} - \frac{B}{6} \;=\; -\frac{B}{12} \;<\; 0 ,
\end{equation}
a contradiction.
\end{proof}

\begin{remark}
The essential point is the bound of Step~2 relating positive-eligible edges to
bridges. Even the weaker estimate $P \le 4B$, obtained by ignoring the degree
condition on the two ends of a bridge, would give total excess at most $0$ and
hence still contradict~\eqref{eq:identity}; the bound $P \le 2B$ supplies the
margin $-B/12$ in~\eqref{eq:final}.
\end{remark}

\begin{remark}\label{rem:bridged}
Theorem~\ref{thm:x2weak} and Proposition~\ref{prop:twoY-weak} are statements
about two-valued graphs, not about bridge-relaxed alternating plane graphs in
general, which need not be bridgeless. Take three copies $H_1,H_2,H_3$ of the
$17$-vertex $(3,4,5)$-alternating plane graph Schneider-17 of
(\cite{AHSSV2015}, Figure~2), each redrawn so that a face incident with a chosen
degree-$5$ vertex $z_i$ is the exterior face, and join $z_2$ to $z_1$ and to
$z_3$ by two bridges drawn in the exterior. Then $\deg z_2 = 7$ and
$\deg z_1 = \deg z_3 = 6$, while their neighbours keep degree $3$ or $4$, so
(c) still holds. The merged exterior face has size
$|F_1|+|F_2|+|F_3|+4 \ge 13$, which differs from every interior face size, so (d)
holds between distinct adjacent faces. The result is a bridge-relaxed
alternating plane graph with bridges. It has neither exactly two degrees nor
exactly two face sizes, and so is no counterexample above.
\end{remark}

\section*{Disclosure of AI use}

Large language models were used in this work and the extent is material. They
assisted with exploration, with adversarial review, and with drafting the
manuscript; the idea of removing the bridge hypothesis, carried out in
\S\ref{sec:c2}, was proposed by an AI reviewer. The author re-derived every
argument independently, has verified the mathematical content, and is answerable
for it and for this disclosure. No AI system is an author.

\section*{Computational support}

The proofs above are self-contained and depend on no computation. A public
repository accompanying this work contains exact-arithmetic checks of the
inequalities used here and machine checks of explicit alternating plane graphs
against Definition~\ref{def:apg}, including the boundary-walk convention of
Convention~\ref{conv:c1}. Those checks are advisory and cover examples, not the
universal statements proved here; no claim of formal verification is made, and a
referee need not run them.

\begin{thebibliography}{9}

\bibitem{AHSSV2015}
I.~Alth\"ofer, J.~K. Haugland, K.~Scherer, F.~Schneider and N.~Van Cleemput,
\emph{Alternating plane graphs}, Ars Math. Contemp. \textbf{8} (2015),
337--363.
\href{https://doi.org/10.26493/1855-3974.584.09a}{doi:10.26493/1855-3974.584.09a}.

\end{thebibliography}

\end{document}
